\chapter{Uveitis}
\index{Uveitis}

\section*{Definition and Overview}
Uveitis is a broad medical term encompassing a complex group of intraocular inflammatory diseases that primarily affect the uvea, or uveal tract—the middle, highly vascular layer of the eye. The uveal tract consists of three main contiguous structures: the iris (the colored part of the eye surrounding the pupil), the ciliary body (the ring-shaped muscle that produces aqueous humor and controls lens accommodation), and the choroid (the vascular layer lining the posterior eye that provides oxygen and nourishment to the outer layers of the retina). While the primary site of inflammation is the uvea, the inflammatory process often extends to adjacent intraocular tissues, including the retina, optic nerve, vitreous humor, sclera, and cornea, potentially leading to profound structural damage and permanent vision impairment.

Anatomically and clinically, uveitis is classified into four primary categories based on the Standardization of Uveitis Nomenclature (SUN) criteria:
\begin{itemize}
    \item \textbf{Anterior Uveitis}: The most common form, where inflammation is localized predominantly in the anterior chamber. It includes iritis (inflammation of the iris) and iridocyclitis (inflammation of both the iris and the ciliary body).
    \item \textbf{Intermediate Uveitis}: Characterized by primary inflammation in the vitreous humor and the pars plana (the posterior portion of the ciliary body). It is often referred to as pars planitis when it is idiopathic and presents with characteristic inflammatory deposits (snowbanks and snowballs) in the vitreous.
    \item \textbf{Posterior Uveitis}: Inflammation is centered in the posterior segment of the eye, affecting the retina (retinitis), the choroid (choroiditis), or both (chorioretinitis/retinochoroiditis).
    \item \textbf{Panuveitis}: A severe, widespread form where inflammation is observed simultaneously across all three anatomical regions (anterior, intermediate, and posterior segments).
\end{itemize}

Understanding uveitis is of paramount clinical significance because it represents a leading cause of preventable blindness worldwide. In developed nations, uveitis is responsible for approximately 10\% to 15\% of cases of legal blindness. Unlike other major vision-threatening eye conditions such as age-related macular degeneration or cataracts, which primarily affect older demographics, uveitis disproportionately affects young and middle-aged adults during their most productive working years. Consequently, the socioeconomic burden of the disease is exceptionally high, resulting in substantial disability, loss of productivity, and high healthcare utilization. Furthermore, uveitis is frequently the ocular manifestation of an underlying, undiagnosed systemic autoimmune or infectious disease, making its recognition a critical diagnostic window into a patient's overall health.

\section*{Epidemiology}
The epidemiology of uveitis is highly diverse, reflecting its multifactorial etiology and wide geographic variations. Globally, the estimated incidence of uveitis ranges from 17 to 52 cases per 100,000 person-years, while the prevalence is estimated between 58 and 115 cases per 100,000 individuals. Although uveitis can occur at any age, it demonstrates a bimodal age distribution, with a primary peak in adults aged 20 to 50 years and a secondary, smaller peak in older adults. Pediatric uveitis is less common, accounting for approximately 5\% to 10\% of all cases, and is strongly associated with juvenile idiopathic arthritis (JIA).

Demographic and geographic patterns are heavily influenced by genetic factors, particularly Human Leukocyte Antigen (HLA) associations, and local exposures to infectious agents:
\begin{itemize}
    \item \textbf{Anterior Uveitis}: This form accounts for approximately 50\% to 90\% of all cases in outpatient clinics. In Western populations, a significant portion of these cases are associated with the HLA-B27 allele. HLA-B27-associated acute anterior uveitis is more prevalent in young males and is frequently linked to seronegative spondyloarthropathies.
    \item \textbf{Beh\c{c}et's Disease}: A systemic vasculitis that causes severe panuveitis, Beh\c{c}et's disease is highly prevalent along the historic ``Silk Road,'' extending from the Mediterranean basin through the Middle East to East Asia. It exhibits a strong association with the HLA-B51 allele.
    \item \textbf{Vogt-Koyanagi-Harada (VKH) Syndrome}: A systemic autoimmune disorder targeting melanocytes, VKH syndrome is a common cause of panuveitis in Asian, Middle Eastern, Native American, and Hispanic populations, with a female predisposition and links to HLA-DR4.
    \item \textbf{Sarcoidosis-Associated Uveitis}: This manifestation is particularly prevalent among African Americans and Scandinavians, showing a female predominance and bimodal age onset.
    \item \textbf{Birdshot Chorioretinopathy}: A chronic posterior uveitis, birdshot chorioretinopathy occurs almost exclusively in individuals of Northern European descent who carry the HLA-A29 allele, representing one of the strongest known genetic associations in medicine.
\end{itemize}

\section*{Aetiology and Pathophysiology}
The aetiology of uveitis is complex and can be broadly categorized into infectious, non-infectious (autoimmune), masquerade syndromes, and idiopathic forms. The pathophysiology involves a breakdown of the blood-ocular barrier, which normally maintains the immune-privileged status of the eye. This barrier consists of tight junctions between the vascular endothelial cells of the retina and iris, as well as the epithelial cells of the ciliary body and retinal pigment epithelium (RPE). When this barrier is disrupted, circulating leukocytes, proteins, and inflammatory cytokines gain access to the intraocular compartments, initiating an inflammatory cascade.

Key causes and mechanisms include:
\begin{itemize}
    \item \textbf{Autoimmune and Autoinflammatory Pathways}: Non-infectious uveitis is primarily a T-cell-mediated autoimmune response. Sensitized CD4+ T-helper cells (specifically Th1 and Th17 lineages) recognize intraocular self-antigens, such as retinal soluble antigen (S-Ag) or interphotoreceptor retinoid-binding protein (IRBP). These T-cells cross the blood-retinal barrier and release pro-inflammatory cytokines, including interferon-gamma (IFN-gamma), interleukin-2 (IL-2), and interleukin-17 (IL-17). These cytokines recruit and activate secondary effector cells (macrophages and neutrophils), which release reactive oxygen species and nitric oxide, causing tissue damage.
    \item \textbf{HLA-B27 and Molecular Mimicry}: The HLA-B27 molecule, a Class I MHC molecule, plays a critical role in antigen presentation. The pathophysiology of HLA-B27-associated uveitis is believed to involve molecular mimicry, where an immune response initiated against a bacterial pathogen (such as \textit{Klebsiella}, \textit{Yersinia}, \textit{Salmonella}, or \textit{Shigella}) in the gut or urinary tract cross-reacts with self-peptides in the uveal tissue, triggering acute, recurrent inflammation.
    \item \textbf{Viral Infections}: Viruses are major causes of infectious uveitis. Herpes Simplex Virus (HSV) and Varicella Zoster Virus (VZV) typically cause unilateral acute anterior uveitis, often associated with sectorial iris atrophy and elevated intraocular pressure due to trabeculitis. Cytomegalovirus (CMV) is a primary cause of necrotizing retinitis in immunocompromised patients (such as those with HIV/AIDS) but can also cause chronic anterior uveitis in immunocompetent individuals. Rubella virus is strongly associated with Fuchs Uveitis Syndrome (FUS).
    \item \textbf{Bacterial Infections}: \textit{Mycobacterium tuberculosis} (Tuberculosis) and \textit{Treponema pallidum} (Syphilis) are classic bacterial pathogens. Syphilis is known as the ``great masquerader'' because it can present as any anatomical form of uveitis, accompanied by retinitis, vasculitis, or optic neuritis. Lyme disease (\textit{Borrelia burgdorferi}) and cat-scratch disease (\textit{Bartonella henselae}) are also important bacterial causes.
    \item \textbf{Parasitic and Fungal Infections}: \textit{Toxoplasma gondii} is the most common cause of infectious posterior uveitis globally, characterized by a focal necrotizing retinochoroiditis adjacent to an old, pigmented scar. Other parasitic causes include \textit{Toxocara canis}. Fungal uveitis, such as that caused by \textit{Candida} or \textit{Aspergillus}, is typically seen in intravenous drug users, immunocompromised patients, or post-surgical cases (endophthalmitis). Presumed Ocular Histoplasmosis Syndrome (POHS) is associated with \textit{Histoplasma capsulatum}.
    \item \textbf{Masquerade Syndromes}: These are non-inflammatory conditions that mimic uveitis. Neoplastic masquerades include intraocular lymphoma (vitreoretinal lymphoma), leukemia, retinoblastoma (in pediatric patients), and uveal melanoma. Non-neoplastic masquerades include intraocular foreign bodies, retinal detachment, and pigment dispersion syndrome.
    \item \textbf{Idiopathic Uveitis}: In approximately 30\% to 50\% of cases, particularly in anterior uveitis, no specific infectious cause or systemic disease association can be identified despite an exhaustive diagnostic workup. These cases are classified as idiopathic and are managed empirically.
\end{itemize}

\section*{Clinical Features}
The clinical presentation of uveitis varies widely depending on the anatomical location, duration (acute vs. chronic), and course of the inflammation (recurrent vs. persistent). Acute uveitis typically presents abruptly and lasts less than three months, whereas chronic uveitis is persistent, lasting more than three months, and often has a gradual, asymptomatic onset.

Key clinical symptoms and signs include:
\begin{itemize}
    \item \textbf{Pain}: Dull, aching, deep-seated eye pain is a hallmark of acute anterior uveitis. The pain is caused by spasms of the ciliary muscle and sphincter pupillae. It often radiates to the temple, forehead, or periorbital region. Pain is characteristically absent in intermediate, posterior, and chronic anterior uveitis.
    \item \textbf{Photophobia}: Extreme sensitivity to light, particularly consensual photophobia (where shining light into the unaffected eye elicits pain in the affected eye), is highly indicative of anterior uveitis.
    \item \textbf{Redness}: A distinctive pattern of injection known as ciliary flush or circumcorneal injection (a violaceous ring of dilated deep episcleral vessels surrounding the limbus) is common in acute anterior uveitis, contrasting with the diffuse conjunctival injection of conjunctivitis.
    \item \textbf{Visual Disturbances}: Blurred vision, haze, or a mild-to-severe drop in visual acuity can occur. In intermediate and posterior uveitis, painless vision loss is the primary presenting symptom, often accompanied by metamorphopsia (distorted vision) or scotomas (blind spots).
    \item \textbf{Floaters}: The perception of spots, webs, or dust-like particles moving across the visual field is a prominent symptom of intermediate and posterior uveitis, resulting from inflammatory cells and protein debris casting shadows on the retina from the vitreous cavity.
    \item \textbf{Keratic Precipitates (KPs)}: These are inflammatory cell deposits on the corneal endothelium. These can be ``fine'' (composed of neutrophils and lymphocytes, typical of acute non-granulomatous uveitis) or large, greasy ``mutton-fat'' KPs (composed of macrophages and epithelioid cells, indicative of granulomatous uveitis, such as in Sarcoidosis or Tuberculosis).
    \item \textbf{Aqueous Flare and Cells}: Visualized on slit-lamp examination. Cells represent individual floating inflammatory cells in the anterior chamber, graded from 0 to 4+. Flare represents protein leakage from inflamed iris vessels, causing a light-scattering effect similar to a projector beam in a dusty room.
    \item \textbf{Hypopyon}: A sterile accumulation of white blood cells pooling at the bottom of the anterior chamber, forming a visible horizontal fluid line. This is commonly seen in HLA-B27 uveitis and Beh\c{c}et's disease.
    \item \textbf{Pupillary Changes}: The pupil of the affected eye is often miotic (constricted) due to sphincter muscle spasm. If posterior synechiae (adhesions between the iris and the anterior lens capsule) form, the pupil may appear irregular or fixed, and fail to dilate properly.
    \item \textbf{Vitreous Haze and Cells}: In intermediate and posterior uveitis, inflammatory cells in the vitreous cavity lead to vitreous haze, which can be graded clinically. Aggregations of inflammatory cells can form ``snowballs'' in the vitreous or deposit on the pars plana as ``snowbanks,'' particularly in the inferior periphery.
    \item \textbf{Chorioretinal Lesions}: Posterior uveitis exhibits active white or yellow-white lesions in the retina or choroid, representing foci of necrotizing retinitis or active choroiditis. Old lesions appear as well-circumscribed areas of chorioretinal atrophy with hyperpigmented borders.
\end{itemize}

\section*{Differential Diagnosis}
Because the symptoms of uveitis (such as red eye, pain, and vision loss) overlap with numerous other ophthalmic and systemic conditions, a broad differential diagnosis must be considered. Misdiagnosis can lead to inappropriate treatment, such as administering corticosteroids in the presence of an active bacterial ulcer, which can have catastrophic consequences.

The differential diagnosis includes:
\begin{itemize}
    \item \textbf{Acute Conjunctivitis}: Unlike uveitis, conjunctivitis is characterized by diffuse conjunctival injection (more intense toward the fornices rather than the limbus), discharge (watery, mucoid, or purulent), itching, and the absence of photophobia, deep pain, ciliary flush, or anterior chamber cells.
    \item \textbf{Corneal Abrasion or Ulcer}: Epithelial defects or infectious keratitis present with pain, redness, and photophobia. Slit-lamp examination with fluorescein staining will reveal corneal epithelial loss or an active stromal infiltrate, which are absent in isolated uveitis.
    \item \textbf{Acute Angle-Closure Glaucoma}: Represents an ophthalmic emergency presenting with sudden, severe unilateral eye pain, headache, nausea, vomiting, blurred vision with halos around lights, and a red eye. Physical findings include a cloudy cornea, a mid-dilated, non-reactive pupil, and extremely elevated intraocular pressure, which distinguishes it from the typically low or normal pressure of acute anterior uveitis.
    \item \textbf{Episcleritis and Scleritis}: Episcleritis is a benign, localized, sectorial redness of the episcleral vessels that blanches with topical phenylephrine and causes mild discomfort. Scleritis is a severe, destructive inflammatory disease presenting with intense, boring eye pain that awakens the patient from sleep, a deep violaceous hue of the sclera, and scleral thinning.
    \item \textbf{Infectious Endophthalmitis}: A devastating bacterial or fungal infection of the intraocular cavities, usually occurring after intraocular surgery, trauma, or hematogenous seeding. It presents with rapid-onset pain, loss of vision, eyelid edema, chemosis, dense hypopyon, and loss of the red reflex, requiring emergent intravitreal antibiotics.
    \item \textbf{Intraocular Lymphoma}: A primary vitreoretinal lymphoma that acts as a classic masquerade syndrome. It must be suspected in older adults presenting with persistent intermediate or posterior uveitis, vitreous cells, and subretinal yellowish infiltrates that do not respond to standard corticosteroid therapy.
    \item \textbf{Retinal Detachment}: Presents with floaters and flashes of light (photopsias) followed by a dark curtain-like shadow over the visual field. It is painless and does not present with anterior chamber inflammation or ciliary flush, though mild vitreous cells can sometimes be present.
\end{itemize}

\section*{When to Seek Medical Care}
Uveitis is a potentially sight-threatening condition that requires prompt professional evaluation. Any delay in diagnosis and treatment can result in irreversible structural damage to the eye, such as permanent scarring of the retina, optic nerve damage, or chronic glaucoma.

Patients must seek immediate medical attention from an optometrist, ophthalmologist, or emergency department if they experience any of the following symptoms:
\begin{itemize}
    \item Sudden, significant, or unexplained reduction in vision or loss of visual acuity.
    \item Severe, persistent, or worsening eye pain, particularly if it radiates to the temple or forehead.
    \item Extreme sensitivity to light (photophobia) that makes it difficult to open the eyes in normal lighting.
    \item Deep, localized redness around the iris (ciliary flush) rather than a generalized pink appearance of the conjunctiva.
    \item The sudden appearance of numerous new floaters, flashes of light, or a dark shadow over the field of vision.
    \item An irregular or misshapen pupil, or a visible layer of pus pooling at the bottom of the colored part of the eye (hypopyon).
\end{itemize}

For life-threatening symptoms associated with systemic conditions (such as a severe, sudden headache accompanied by a stiff neck, high fever, confusion, or focal neurological deficits, which may indicate systemic meningitis, Beh\c{c}et's vasculitis flare, or Vogt-Koyanagi-Harada syndrome), patients must immediately call emergency services. The emergency number is \textbf{000} in many countries, \textbf{911} in the United States and Canada, \textbf{999} in the United Kingdom, and \textbf{112} in the European Union.

\section*{Diagnosis and Investigations}
The diagnosis of uveitis involves a comprehensive ophthalmic evaluation combined with a targeted systemic workup. Because uveitis is a manifestation of numerous systemic diseases, the diagnostician must act as a medical detective, using clinical clues from the eye exam to guide laboratory and imaging investigations, rather than ordering a non-specific screening panel.

Diagnostic procedures and investigations include:
\begin{itemize}
    \item \textbf{Ophthalmic Examination}:
    \begin{itemize}
        \item \textbf{Visual Acuity Testing}: Measures the baseline level of vision impairment.
        \item \textbf{Slit-Lamp Biomicroscopy}: The primary tool to evaluate the anterior segment. It allows the clinician to identify and grade keratic precipitates, anterior chamber cells and flare, posterior synechiae, iris nodules (Koeppe and Busacca nodules), and iris atrophy.
        \item \textbf{Applanation Tonometry}: Measures intraocular pressure. Pressure is often low in acute uveitis due to ciliary body hyposecretion, but can be elevated (hypertensive uveitis) in herpetic or cytomegalovirus infections, or as a side effect of corticosteroid therapy.
        \item \textbf{Dilated Fundus Examination}: Utilizing indirect ophthalmoscopy to view the vitreous cavity, retina, choroid, blood vessels, and optic nerve head to detect vitreous cells, snowbanks, retinitis, choroiditis, vasculitis, macular edema, or optic disc swelling.
    \end{itemize}
    \item \textbf{Ocular Imaging}:
    \begin{itemize}
        \item \textbf{Optical Coherence Tomography (OCT)}: A non-invasive imaging modality that provides high-resolution, cross-sectional images of the retina. It is the gold standard for diagnosing and monitoring cystoid macular edema (CME) and subretinal fluid.
        \item \textbf{Fluorescein Angiography (FA)}: Involves injecting a fluorescent dye intravenously to visualize the retinal circulation. It is essential for detecting retinal vasculitis (demonstrated by vascular leakage and ferning), capillary non-perfusion, and choroidal neovascularization.
        \item \textbf{Indocyanine Green Angiography (ICGA)}: Utilizes a dye that fluoresces in the near-infrared spectrum, allowing visualization of the deeper choroidal vasculature. It is particularly useful in white dot syndromes, VKH, and sympathetic ophthalmia.
        \item \textbf{B-Scan Ultrasonography}: Used when dense media (such as a mature cataract or severe vitreous hemorrhage) prevents direct visualization of the posterior segment, helping to rule out retinal detachment or intraocular tumors.
    \end{itemize}
    \item \textbf{Laboratory and Systemic Investigations}:
    \begin{itemize}
        \item \textbf{Human Leukocyte Antigen (HLA) Typing}: Specifically HLA-B27 for patients with recurrent, acute unilateral anterior uveitis, or HLA-A29 if birdshot chorioretinopathy is suspected.
        \item \textbf{Syphilis Serology}: Screening with non-treponemal tests (RPR or VDRL) and confirmation with treponemal-specific tests (FTA-ABS or TPPA) is mandatory in almost all cases of unexplained uveitis, as syphilis can mimic any presentation.
        \item \textbf{Tuberculosis Screening}: Interferon-Gamma Release Assays (IGRA, e.g., QuantiFERON-TB Gold) or a purified protein derivative (PPD) Mantoux skin test.
        \item \textbf{Sarcoidosis Workup}: Serum Angiotensin-Converting Enzyme (ACE) and serum lysozyme levels, combined with a Chest X-ray or high-resolution Chest CT scan to identify bilateral hilar lung changes.
        \item \textbf{Autoantibody Screening}: Antinuclear Antibodies (ANA) are critical in pediatric patients, as ANA-positivity in children with juvenile idiopathic arthritis (JIA) carries a high risk of developing asymptomatic, sight-threatening chronic anterior uveitis. Anti-Neutrophil Cytoplasmic Antibodies (ANCA) are ordered if granulomatosis with polyangiitis is suspected.
    \end{itemize}
    \item \textbf{Invasive Diagnostic Procedures}:
    \begin{itemize}
        \item \textbf{Anterior Chamber Paracentesis or Vitreous Aspiration}: Obtaining aqueous or vitreous humor samples for Polymerase Chain Reaction (PCR) testing (highly sensitive for viral and parasitic DNA, such as HSV, VZV, CMV, and \textit{Toxoplasma}), cultures (for atypical bacteria or fungi), or cytology (to rule out vitreoretinal lymphoma).
    \end{itemize}
\end{itemize}

\section*{Management}
The primary goals of uveitis management are to suppress ocular inflammation, alleviate symptoms (such as pain and photophobia), prevent structural complications, treat any underlying infection, and preserve visual function. Management must be tailored to the anatomical location, severity, and etiology of the disease.

Therapeutic options include:
\begin{itemize}
    \item \textbf{Mydriatics and Cycloplegics}: Topical agents such as cyclopentolate 1\%, homatropine 2\%, or atropine 1\% are crucial in acute anterior uveitis. By paralyzing the ciliary muscle and sphincter pupillae, they relieve pain from ciliary spasm. They also dilate the pupil, preventing the formation of posterior synechiae or breaking recently formed adhesions.
    \item \textbf{Corticosteroids}: The cornerstone of non-infectious uveitis therapy.
    \begin{itemize}
        \item \textbf{Topical Corticosteroids}: Formulations like prednisolone acetate 1\% or difluprednate 0.05\% are the standard of care for anterior uveitis. Dosing frequency can range from hourly in severe acute cases to a slow, gradual taper over several weeks or months.
        \item \textbf{Local Injections}: Periocular (sub-Tenon's or transeptal) or intravitreal injections of triamcinolone acetonide are highly effective for unilateral intermediate or posterior uveitis, especially when complicated by macular edema. Intravitreal corticosteroid implants (such as dexamethasone or fluocinolone acetonide) provide sustained drug release over several months to years.
        \item \textbf{Systemic Corticosteroids}: Oral prednisone is indicated for bilateral, severe, or sight-threatening intermediate, posterior, or panuveitis, or when uveitis is associated with systemic disease. A typical starting dose is 1 mg/kg/day, followed by a slow, supervised taper to avoid rebound inflammation.
    \end{itemize}
    \item \textbf{Systemic Immunomodulatory Therapy (IMT)}: Indicated when uveitis is sight-threatening and refractory to corticosteroids, when corticosteroid-sparing therapy is required due to side effects, or in specific diseases where early IMT is the standard of care (such as Beh\c{c}et's disease or VKH).
    \begin{itemize}
        \item \textbf{Antimetabolites}: Methotrexate, Mycophenolate Mofetil (MMF), and Azathioprine are first-line immunomodulators. They inhibit lymphocyte proliferation.
        \item \textbf{T-cell Inhibitors}: Cyclosporine and Tacrolimus inhibit calcineurin pathways, suppressing T-cell activation.
        \item \textbf{Biologic Response Modifiers}: Monoclonal antibodies targeting Tumor Necrosis Factor-alpha (TNF-alpha), such as Adalimumab and Infliximab, have revolutionized the management of severe, non-infectious uveitis. Adalimumab is the first FDA-approved non-corticosteroid systemic therapy for non-infectious intermediate, posterior, and panuveitis.
    \end{itemize}
    \item \textbf{Antimicrobial Therapy}: Essential for infectious uveitis. Corticosteroids are used with caution and always in combination with specific anti-infectives to control the inflammatory response without promoting pathogen replication.
    \begin{itemize}
        \item \textbf{Toxoplasmosis}: Treated with a combination of pyrimethamine, sulfadiazine, and folinic acid (triple therapy), or alternatively with oral clindamycin, trimethoprim/sulfamethoxazole, or intravitreal injections of clindamycin and dexamethasone.
        \item \textbf{Herpetic Infections}: Oral or intravenous acyclovir, valacyclovir, or famciclovir for HSV and VZV. CMV retinitis requires induction and maintenance therapy with ganciclovir (intravenous or intravitreal) or oral valganciclovir.
        \item \textbf{Tuberculosis and Syphilis}: Require standard multidrug antitubercular therapy (Rifampin, Isoniazid, Pyrazinamide, Ethambutol) and intravenous Penicillin G, respectively, managed in collaboration with infectious disease specialists.
    \end{itemize}
    \item \textbf{Surgical Management}:
    \begin{itemize}
        \item \textbf{Pars Plana Vitrectomy (PPV)}: Performed to remove vitreous opacities, relieve vitreomacular traction, peel epiretinal membranes, or diagnostic biopsy.
        \item \textbf{Cataract Extraction}: Indicated when uveitic cataracts impair vision or prevent visualization of the fundus. Surgery must only be performed after the eye has been completely quiet (zero cells) for at least three consecutive months.
        \item \textbf{Glaucoma Surgery}: Trabeculectomy or glaucoma drainage device (tube shunt) implantation when intraocular pressure cannot be controlled with topical medications.
    \end{itemize}
\end{itemize}

\section*{Complications}
Chronic or recurrent intraocular inflammation, as well as the long-term therapies used to treat it, can lead to severe ocular complications. Managing these complications is often as challenging as treating the primary inflammation.

Potential complications include:
\begin{itemize}
    \item \textbf{Cystoid Macular Edema (CME)}: The accumulation of fluid in the outer plexiform layer of the central retina (macula), leading to cystic spaces. CME is the leading cause of permanent vision loss in patients with uveitis, particularly intermediate, posterior, and panuveitis.
    \item \textbf{Secondary Glaucoma}: Elevated intraocular pressure that causes optic nerve damage. It can occur via:
    \begin{itemize}
        \item \textbf{Open-Angle Glaucoma}: Caused by inflammatory debris, cells, or proteins clogging the trabecular meshwork, or due to corticosteroid-induced changes in trabecular outflow (steroid-response glaucoma).
        \item \textbf{Angle-Closure Glaucoma}: Caused by pupillary block due to 360-degree posterior synechiae (seclusio pupillae), which prevents aqueous humor from moving from the posterior to the anterior chamber, causing the iris to bow forward (iris bombé) and close the iridocorneal angle.
    \end{itemize}
    \item \textbf{Cataract Formation}: Posterior subcapsular cataracts are highly common. They develop due to the direct effect of chronic intraocular inflammation on lens metabolism, or as a well-documented dose- and duration-dependent complication of local or systemic corticosteroid therapy.
    \item \textbf{Band Keratopathy}: The deposition of calcium phosphate salts in the superficial layers of the cornea (Bowman's membrane). It is characteristically seen in chronic pediatric uveitis associated with juvenile idiopathic arthritis (JIA) and can cause ocular irritation and decreased vision.
    \item \textbf{Posterior Synechiae}: Adhesions between the posterior iris and the anterior lens capsule. This can result in a distorted, non-reactive pupil and predispose the patient to angle-closure glaucoma.
    \item \textbf{Vitreous Opacification and Hemorrhage}: Severe vitreous haze or traction on retinal vessels can cause vitreous hemorrhage, severely obstructing vision.
    \item \textbf{Retinal Detachment}: Tractional retinal detachment can occur due to the contraction of inflammatory membranes in the vitreous cavity. Rhegmatogenous retinal detachment can occur if necrotizing retinitis (such as in acute retinal necrosis) causes retinal tears.
    \item \textbf{Choroidal Neovascularization (CNV)}: The growth of abnormal, fragile new blood vessels from the choroid through breaks in Bruch's membrane into the subretinal space. These vessels can leak fluid or bleed, causing rapid, severe central vision loss.
    \item \textbf{Phthisis Bulbi}: The final, end-stage clinical state of an eye that has suffered severe, chronic, uncontrolled inflammation. The ciliary body undergoes atrophy, leading to permanent hypotony (extremely low intraocular pressure), structural collapse, shrinkage, and complete loss of function.
\end{itemize}

\section*{Prognosis and Outcomes}
The prognosis for patients with uveitis is highly variable and depends on several factors, including the anatomical classification, the underlying etiology, the severity of the inflammation at presentation, and the promptness of medical intervention.

For patients with acute anterior uveitis, the prognosis is generally excellent. Most episodes resolve completely within days to a few weeks with appropriate topical corticosteroid and cycloplegic therapy, without leaving permanent visual deficits. However, recurrence is common, especially in HLA-B27-positive individuals, and recurrent episodes require prompt treatment to prevent cumulative damage.

Conversely, chronic anterior, intermediate, posterior, and panuveitis carry a more guarded prognosis. These forms of the disease are prone to persistent, low-grade inflammation that gradually damages ocular structures over years. The prognosis is significantly worse for patients who develop complications such as cystoid macular edema, glaucoma, or choroidal neovascularization. Delayed diagnosis or inadequate treatment increases the risk of permanent, legal blindness.

The prognosis is also heavily influenced by the underlying cause:
\begin{itemize}
    \item \textbf{Infectious Uveitis}: If diagnosed early and treated with targeted antimicrobials (e.g., in toxoplasmosis or syphilis), the infection can be eradicated, though permanent scars may remain. Viral uveitis (HSV/VZV) tends to be recurrent and requires long-term monitoring.
    \item \textbf{Systemic Autoimmune Disease}: Uveitis associated with conditions like Beh\c{c}et's disease, sarcoidosis, or VKH syndrome requires close, lifelong multidisciplinary management. While modern biologic therapies (like adalimumab) have significantly improved visual outcomes and drug-free remission rates, these patients remain at risk for chronic relapses.
    \item \textbf{Pediatric Uveitis}: JIA-associated uveitis is notoriously asymptomatic (a ``quiet'' red eye) but highly destructive, often presenting only after complications like cataracts or band keratopathy have already developed. Regular screening is essential for a favorable long-term prognosis.
\end{itemize}

\section*{Prevention and Self-Care}
While many forms of autoimmune and idiopathic uveitis cannot be prevented, there are practical steps patients can take to reduce the risk of infectious uveitis, minimize the frequency and severity of recurrences, and manage the disease effectively.

Prevention and self-care strategies include:
\begin{itemize}
    \item \textbf{Smoking Cessation}: Cigarette smoking is a strong, independent risk factor for both the development of uveitis and increased disease activity. Smoking is associated with a higher frequency of flares, a greater risk of macular edema, and a poor response to treatment. Quitting smoking is one of the most effective lifestyle changes a patient can make.
    \item \textbf{Dietary and Lifestyle Hygiene}: To prevent infectious causes of posterior uveitis:
    \begin{itemize}
        \item Cook meat thoroughly to safe internal temperatures to prevent infection with \textit{Toxoplasma gondii}.
        \item Wash hands thoroughly after handling raw meat, gardening, or cleaning a cat's litter box (cats are the definitive host for \textit{Toxoplasma}).
        \item Practice safe sex (such as consistent condom use) to prevent the transmission of syphilis and HIV, both of which are major causes of infectious uveitis.
    \end{itemize}
    \item \textbf{Strict Medication Adherence}: Patients must follow their prescribed eye drop and oral medication schedules precisely. Corticosteroid eye drops must never be stopped abruptly, as this can trigger a severe, rebound inflammatory flare. Tapering must always be done gradually under the direct supervision of an ophthalmologist.
    \item \textbf{Home Monitoring}: Patients with chronic or posterior uveitis should perform regular self-monitoring using an Amsler grid to detect early signs of macular edema or choroidal neovascularization, which present as wavy or distorted lines (metamorphopsia). Any changes must be reported immediately.
    \item \textbf{Protecting the Eyes}: Wearing sunglasses with UV protection can help reduce discomfort from severe photophobia during acute flares.
    \item \textbf{Adhering to Routine Screening}: Children diagnosed with juvenile idiopathic arthritis (JIA) must undergo regular, scheduled slit-lamp screenings by an ophthalmologist, even if they have no visual symptoms, as early detection is key to preventing permanent blindness.
\end{itemize}

\section*{Sources and Further Reading}
For authoritative, up-to-date, and evidence-based information regarding uveitis and its management, the following resources are recommended:
\begin{itemize}
    \item \textbf{National Eye Institute (NEI)}: Provides comprehensive patient education materials, research updates, and clinical trials information. \\
    URL: https://www.nei.nih.gov/learn-about-eye-health/eye-conditions-and-diseases/uveitis
    \item \textbf{American Academy of Ophthalmology (AAO)}: Offers detailed clinical summaries, patient guides, and expert-reviewed articles on the diagnosis and treatment of all forms of uveitis. \\
    URL: https://www.aao.org/eye-health/diseases/what-is-uveitis
    \item \textbf{Mayo Clinic}: A world-renowned medical institution providing clear information on symptoms, causes, diagnostic tests, and treatment options. \\
    URL: https://www.mayoclinic.org/diseases-conditions/uveitis/symptoms-causes/syc-20378734
    \item \textbf{Cleveland Clinic}: Offers clinical guides on the classification of uveitis, its relationship to systemic diseases, and the role of immunosuppressive therapy. \\
    URL: https://my.clevelandclinic.org/health/diseases/17614-uveitis
    \item \textbf{Uveitis SG (Standardization of Uveitis Nomenclature Consortium)}: Access to academic publications, diagnostic criteria guidelines, and classification tools developed by the international SUN working group. \\
    URL: https://www.clinics.org.uk/sun-classification
    \item \textbf{National Health Service (NHS)}: The United Kingdom's public health service provides practical patient guidance on recognizing symptoms, treatment pathways, and when to seek emergency eye care. \\
    URL: https://www.nhs.uk/conditions/uveitis/
\end{itemize}